% ==============================================================================
% RECOMENDACIONES
% ==============================================================================
% Solo componente de ciberseguridad. Cada recomendación responde a una
% conclusión de m-conclusiones.tex o a una restricción declarada en
% j-metodologia.tex. La recomendación de validación con una startup se
% incorporó por pedido explícito del estudiante (2026-10-09).
\ifdefined\modoensamble
\else
	\documentclass[11pt, letterpaper]{report}
\usepackage[spanish, es-nodecimaldot, es-noquoting]{babel}
\selectlanguage{spanish}
\usepackage[utf8]{inputenc}
\usepackage[T1]{fontenc}

% ===============================================================================
% El estudiante debe llenar sus datos en esta sección para que la plantilla los 
% auto-importe y genere automáticamente las páginas de portada y de firmas 
% autorizadas.
% ===============================================================================
% Datos del estudiante:
% -------------------------------------------------------------------------------
% Nombre completo
\def \nombreestudiante {Sebastián Huertas Gómez}
% Carné
\def \uvgcarne {22295}
% Facultad
\def \uvgfacultad {Ingeniería}
% Carrera
\def \uvgcarrera {Ingeniería en Ciencia de la Computación y Tecnologías de la Información}

% Datos del trabajo:
% -------------------------------------------------------------------------------
% Título completo (tomado literalmente de protocolo_sebastian_v1 (3).pdf, sección 1)
\def \titulotesis {logSguarDian: Desarrollo de una librería para la detección de amenazas de seguridad en Endpoints de aplicaciones}
% Mes y año de entrega (confirmado por el estudiante: noviembre 2026)
\def \mesentrega {noviembre}
\def \anoentrega {2026}
% Asesor (según protocolo_sebastian_v1 (3).pdf, portada)
\def \nombreasesor {Ing. Erick Francisco Marroquín Rodríguez}
% El catedrático del curso (Ing. Bacilio Alexander Bolaños Lima) figura como
% primer examinador en la hoja de firmas, definido más abajo.

% Datos del tribunal examinador:
% -------------------------------------------------------------------------------
% Nombre del primer examinador (segunda línea de la hoja de firmas)
\def \nombreprimerex {Ing. Ma. Bacilio Alexander Bolaños Lima}

% Capítulos pre-definidos
% -------------------------------------------------------------------------------
% Comentar las líneas de las secciones que desean omitirse, por defecto se
% se incluyen todas.
% Se descomentan conforme se redacta cada capítulo (todavía sin archivo de
% contenido: prefacio, alcance, anexos, glosario).
\def \CAPprefacio {Prefacio}
\def \CAPantecedentes {Antecedentes}
\def \CAPalcance {Alcance}
\def \CAPmetodologia {Metodología}
\def \CAPresultados {Resultados}
\def \CAPdiscusion {Análisis de resultados}
\def \CAPanexos {Anexos}
\def \CAPglosario {Glosario}

% Formato y estilo de la plantilla
% -------------------------------------------------------------------------------
% Modo impresión: Puede des-comentar la siguiente línea para generar un documento pdf sin la portada, para cuando se desee imprimir el documento para encuadernación
% NOTA: la plantilla original traía esta línea activa por defecto, lo cual
% suprime la portada verde (main.tex hace \let\CAPportada\undefined cuando
% \printver está definido). Se comenta para que la versión digital, con
% portada, sea la que se genera por defecto.
%\def \printver {Versión del documento para impresión}

% Portada: Puede cambiarse la imagen en la portada al cambiar el nombre del 
% archivo siguiente. NOTA: debe tener la suficiente resolución para cubrir el área
% designada
\def \imagenportada {plantilla/portadacit.jpg}

% Referencias: Puede des-comentar la siguiente línea para utilizar el formato de referencias APA
%\def \usarAPA {Usar formato APA}

% Párrafo: Puede comentar la siguiente línea si desea emplear un formato de 
% párrafo distinto al establecido por defecto
\def \parpordefecto {Formato de párrafo por defecto}

% Capítulos y secciones: Puede des-comentar la siguiente línea para establecer el
% formato de los capítulos y secciones bajo el estándar original de UVG para
% trabajos de graduación. Este incluye: capítulos con numeración romana, secciones
% con letras mayúsculas, sub-secciones con números y sub-sub-secciones con letras
% minúsculas
%\def \capsecuvg {Formato UVG para capítulos y secciones}
% ================================================================================
% En este archivo se colocan opciones adicionales para modificar el formato de la
% plantilla, para emplearse en otros tipos de documentos que no sean trabajos de
% graduación. Si usted está trabajando su tesis, NO modifique este archivo
% ================================================================================
% Capítulos pre-definidos
% --------------------------------------------------------------------------------
% Comentar las líneas de las secciones que desean omitirse, por defecto se 
% se incluyen todas.
\def \CAPportada {Portada}
\def \CAPcaratula {Caratula}
\def \CAPfirmas {Hoja de firmas}
\def \CAPindice {Índice general}
\def \CAPfiguras {Listado de figuras}
\def \CAPcuadros {Listado de cuadros}
\def \CAPresumen {Resumen}
\def \CAPabstract {Resumen}
\def \CAPintroduccion {Introducción}
\def \CAPobjetivos {Objetivos}
\def \CAPjustificacion {Justificación}
\def \CAPmarcoteorico {Marco teórico}
\def \CAPconclusiones {Conclusiones}
\def \CAPrecomendaciones {Recomendaciones}
\def \CAPbibliografia {Bibliografía}

\usepackage{xcolor}
\usepackage{amsfonts}
\usepackage{amsmath}
\usepackage{amssymb}
\usepackage{amsthm}
\usepackage{mathtools}
\usepackage{graphicx}
\usepackage{xfrac}
\usepackage{float}
\usepackage[hypertexnames=false]{hyperref}
\usepackage[font=small]{caption}
\usepackage{subcaption}
% Cuadros: array (columnas L{} alineadas a la izquierda) y longtable (cuadros
% que ocupan más de una página). Se cargan aquí, no en
% 2-paquetes_y_comandos_usuario.tex, para que los capítulos de Metodología y
% Resultados no dependan de ese archivo.
\usepackage{array}
\usepackage{longtable}
\newcolumntype{L}[1]{>{\raggedright\arraybackslash}p{#1}}
\renewcommand{\arraystretch}{1.3}
\setlength{\LTpre}{0.3cm}
\setlength{\LTpost}{0.1cm}
% Margen de estiramiento para que los párrafos con palabras largas no
% generen líneas que sobresalgan del margen.
\setlength{\emergencystretch}{3em}
\usepackage{xpatch}
\usepackage{emptypage}
% htt: permite dividir con guion el texto en \texttt (nombres de archivo,
% clases y campos), que de otro modo sobresale del margen derecho.
\usepackage[htt]{hyphenat}
\usepackage{fancyhdr}
\usepackage[backend=biber, style=ieee]{biblatex}
\ifdefined\usarAPA
    \usepackage[backend=biber, style=apa]{biblatex}
\fi
\addbibresource{m-bibliografia.bib}

\usepackage[percent]{overpic}
\usepackage{chngcntr}

\ifdefined\CAPglosario
	\usepackage[numberedsection]{glossaries}
	\makeglossaries
    \newglossaryentry{rasp}
{
    name={RASP},
    description={Autoprotección de aplicaciones en tiempo de ejecución (\textit{Runtime Application Self-Protection}): mecanismo de seguridad que se ejecuta dentro de la propia aplicación y analiza sus peticiones para detectar y bloquear ataques.}
}

\newglossaryentry{waf}
{
    name={WAF},
    description={Cortafuegos de aplicaciones web (\textit{Web Application Firewall}): dispositivo o servicio perimetral que filtra el tráfico HTTP antes de que llegue a la aplicación.}
}

\newglossaryentry{modsecurity}
{
    name={ModSecurity},
    description={Motor de WAF de código abierto que se usa con reglas externas; en este trabajo se configuró con OWASP CRS.}
}

\newglossaryentry{crs}
{
    name={OWASP CRS},
    description={Conjunto de reglas genéricas de OWASP (\textit{Core Rule Set}) para ModSecurity, con niveles de paranoia que controlan qué tan estrictas son.}
}

\newglossaryentry{owasp}
{
    name={OWASP Top 10},
    description={Lista de las diez categorías de riesgo más críticas para aplicaciones web, publicada por OWASP; este trabajo usa la edición 2021.}
}

\newglossaryentry{mitre}
{
    name={MITRE ATT\&CK},
    sort={MITRE ATT&CK},
    description={Base de conocimiento de tácticas y técnicas de atacantes, usada aquí para asociar cada vector de ataque a una técnica.}
}

\newglossaryentry{cwe}
{
    name={CWE},
    description={Enumeración común de debilidades de software (\textit{Common Weakness Enumeration}): catálogo de tipos de defectos que originan vulnerabilidades.}
}

\newglossaryentry{sqli}
{
    name={SQL Injection},
    description={Ataque que inserta sentencias SQL en una entrada de la aplicación para alterar las consultas que ejecuta la base de datos.}
}

\newglossaryentry{xss}
{
    name={XSS},
    description={\textit{Cross-Site Scripting}: ataque que inyecta código ejecutable en una página que otros usuarios verán.}
}

\newglossaryentry{pathtraversal}
{
    name={Path Traversal},
    description={Ataque que manipula rutas de archivo (por ejemplo con \texttt{../}) para leer archivos fuera del directorio previsto; incluye la inclusión local de archivos (LFI).}
}

\newglossaryentry{cmdi}
{
    name={Command Injection},
    description={Ataque que logra ejecutar comandos del sistema operativo a través de una entrada de la aplicación.}
}

\newglossaryentry{middleware}
{
    name={\textit{middleware}},
    sort={middleware},
    description={Función que se ejecuta entre la recepción de una petición y la lógica de negocio de la aplicación; en Express se registra con \texttt{app.use}.}
}

\newglossaryentry{rf}
{
    name={Random Forest},
    description={Algoritmo supervisado que combina muchos árboles de decisión; en la librería es el único modelo con autoridad para bloquear una petición.}
}

\newglossaryentry{iforest}
{
    name={Isolation Forest},
    description={Algoritmo no supervisado de detección de anomalías que aísla las observaciones atípicas con particiones aleatorias; en la librería solo registra anomalías y nunca bloquea.}
}

\newglossaryentry{onnx}
{
    name={ONNX},
    description={Formato abierto para serializar modelos de aprendizaje automático, que permite ejecutarlos fuera de Python.}
}

\newglossaryentry{workerthreads}
{
    name={\texttt{worker\_threads}},
    sort={worker_threads},
    description={Módulo de Node.js que ejecuta código en hilos independientes del hilo principal; la inferencia de los modelos se ejecuta ahí.}
}

\newglossaryentry{eventloop}
{
    name={\textit{Event Loop}},
    sort={Event Loop},
    description={Bucle de eventos de Node.js que atiende las peticiones en un único hilo; bloquearlo degrada a toda la aplicación.}
}

\newglossaryentry{failopen}
{
    name={\textit{fail-open}},
    sort={fail-open},
    description={Política por la cual, ante un error o un tiempo de espera del módulo de detección, la petición se deja pasar en lugar de bloquearse.}
}

\newglossaryentry{monitor}
{
    name={Modo \texttt{monitor}},
    sort={Modo monitor},
    description={Modo de operación en el que la librería registra las peticiones que bloquearía, pero no las bloquea.}
}

\newglossaryentry{block}
{
    name={Modo \texttt{block}},
    sort={Modo block},
    description={Modo de operación en el que la librería responde HTTP 403 a las peticiones que clasifica como ataque.}
}

\newglossaryentry{pesp}
{
    name={p95},
    description={Percentil 95 del tiempo de respuesta: el valor por debajo del cual queda el 95~\% de las peticiones.}
}

\newglossaryentry{deltap95}
{
    name={$\Delta p95$},
    sort={Delta p95},
    description={Diferencia entre el p95 con el middleware activo y el p95 de la línea base sin él; en este trabajo se mide en milisegundos.}
}

\newglossaryentry{funo}
{
    name={F1},
    description={Media armónica de la precisión y el recall; es la métrica principal de clasificación del trabajo.}
}

\newglossaryentry{auc}
{
    name={AUC-ROC},
    description={Área bajo la curva ROC; mide qué tan bien separa un modelo dos clases sin depender de un umbral.}
}

\newglossaryentry{wilson}
{
    name={Intervalo de Wilson},
    description={Intervalo de confianza para una proporción, adecuado para muestras pequeñas o proporciones cercanas a 0 o 1.}
}

\newglossaryentry{casidup}
{
    name={Casi-duplicado},
    description={Registro muy parecido a otro pero no idéntico; si queda repartido entre entrenamiento y prueba, puede inflar las métricas.}
}

\newglossaryentry{minhash}
{
    name={MinHash},
    description={Técnica para estimar rápidamente la similitud entre conjuntos, usada para detectar casi-duplicados.}
}

\newglossaryentry{fuga}
{
    name={Fuga de datos},
    description={Presencia en el conjunto de evaluación de información disponible durante el entrenamiento, lo que sobreestima el desempeño.}
}

\newglossaryentry{corpuslimpio}
{
    name={Corpus limpio},
    description={Conjunto de payloads de evaluación filtrado para excluir los que coinciden con el entrenamiento; en este trabajo es el de la Ronda~5.}
}

\newglossaryentry{useragent}
{
    name={\texttt{User-Agent}},
    sort={User-Agent},
    description={Cabecera HTTP con la que el cliente se identifica (navegador, cURL, entre otros); su representación en el corpus sesgó al modelo.}
}

\newglossaryentry{falsopositivo}
{
    name={Falso positivo},
    description={Petición legítima que el sistema clasifica o bloquea como ataque.}
}

\newglossaryentry{parmodelos}
{
    name={Par de modelos},
    description={Combinación de un archivo \texttt{rf.onnx} y uno \texttt{if.onnx} identificada por su suma SHA-256 (pares A, B y C).}
}

\newglossaryentry{tarball}
{
    name={\textit{tarball}},
    sort={tarball},
    description={Archivo comprimido \texttt{.tgz} de un paquete de npm, publicado o construido localmente.}
}

\newglossaryentry{kanban}
{
    name={Kanban},
    description={Metodología de gestión basada en un tablero de columnas que representan el estado de las tareas.}
}

\newglossaryentry{artillery}
{
    name={Artillery},
    description={Herramienta de pruebas de carga usada para medir la latencia con y sin el middleware.}
}

\newglossaryentry{sqlmap}
{
    name={sqlmap},
    description={Herramienta de código abierto que automatiza la detección y la explotación de inyecciones SQL; generó los payloads de SQL Injection del corpus limpio.}
}

\newglossaryentry{zap}
{
    name={OWASP ZAP},
    description={Herramienta de OWASP para pruebas de seguridad de aplicaciones web; se usó para confirmar el payload de XSS.}
}

\newglossaryentry{mcp}
{
    name={MCP},
    description={Protocolo de Contexto de Modelo (\textit{Model Context Protocol}): estándar para que un asistente de inteligencia artificial use herramientas externas.}
}

\fi

\usepackage[top=1in, left=1.5in, right=1in, bottom=1in]{geometry}
\usepackage[Sonny]{fncychap}

\graphicspath{ {figuras/} }
\definecolor{uvg-green}{RGB}{17,71,52}
\newcommand{\defaultparformat}[1]{
	{\setlength{\parskip}{2ex}
     \input{#1}}
}
\ifdefined\capsecuvg
	\renewcommand\thechapter{\Roman{chapter}}
    \renewcommand\thesection{\Alph{section}}
	\renewcommand\thesubsection{\arabic{subsection}}
    \renewcommand\thesubsubsection{\alph{subsubection}}
\fi
\counterwithout{figure}{chapter}
\counterwithout{table}{chapter}
\counterwithout{equation}{chapter}

\providecommand{\blankpage}{
\newpage
\thispagestyle{empty}
\mbox{}
\newpage
}

% ==============================================================================
% Paquetes y comandos adicionales del usuario
% ==============================================================================
% TikZ: requerido por las figuras del capítulo de Marco Teórico (pipeline de
% middleware, comparación RASP vs. control perimetral).
\usepackage{tikz}
\usetikzlibrary{arrows.meta, positioning, shapes.geometric, fit, backgrounds}


\pagestyle{headings}
\begin{document}

\ifdefined\parpordefecto
	\setlength{\parskip}{2ex}
\fi

\fi

\chapter{Recomendaciones}
\label{cap:recomendaciones}

Las recomendaciones se derivan de las conclusiones (Capítulo~\ref{cap:conclusiones}), de las
limitaciones (Sección~\ref{sec:dis-limitaciones}) y de las restricciones declaradas en la
Metodología (Sección~\ref{sec:met-restricciones}). Formalizan las sugerencias del capítulo de
Análisis de resultados, que se presentaron allí como indicios. Cada una indica la acción propuesta,
su justificación en la evidencia del trabajo y el resultado que se espera obtener. Se agrupan por
destinatario y por tipo de acción.

\section{Para quienes adopten la librería}

\begin{enumerate}
    \item \textbf{Instalar primero en modo \texttt{monitor}.} \textit{Acción:} instalar la librería
    en modo \texttt{monitor}, revisar durante un período de observación los eventos registrados con
    la CLI (\texttt{attacks list} y \texttt{endpoints top}), ajustar el umbral con el tráfico propio
    de la aplicación y solo después activar el modo \texttt{block}. \textit{Justificación:} con un
    \texttt{User-Agent} de navegador, la librería bloqueó el 92.5~\% del tráfico legítimo de
    dvws-node (Conclusión~\ref{conc:fp}). \textit{Resultado esperado:} conocer la tasa de bloqueos
    indebidos de la propia aplicación antes de afectar a usuarios reales.

    \item \textbf{Usar la librería como capa complementaria detrás de un WAF.}
    \textit{Acción:} desplegarla detrás de un WAF y no como su sustituto. \textit{Justificación:}
    (A) la librería agrega cobertura a un WAF: en la Ronda~5, sobre un corpus filtrado contra el
    entrenamiento, suma 48 de 800 payloads al WAF solo (intervalo de Wilson al 95~\%: 4.6--7.9~\%) y
    bloquea 48 de los 51 ataques que este no bloquea (84.1--98.0~\%); la Ronda~4 lo corrobora con una
    muestra de 40 (37 de 40; 80.1--97.4~\%), pero se midió con el tarball local de la 0.1.0 y el
    ejecutor anterior, y la 3a en otra rama de la aplicación, de modo que no es un resultado de la 0.1.1 (Conclusión~\ref{conc:waf}). (B) La
    librería no sustituye al WAF: en condiciones equiparadas el WAF suma solo 3 de 800 payloads a la
    librería (0.1--1.1~\%), por lo que la cobertura de la Ronda~5 no respalda esta parte; se sostiene
    en que la librería se evaluó solo en cuatro clases de ataque, en que su veredicto depende de un
    voto de árbol y de las cabeceras del despliegue y en los bloqueos indebidos medidos
    (Conclusión~\ref{conc:fp}), sin que se haya medido la tasa de falsos positivos del WAF.
    \textit{Resultado esperado:} cobertura agregada no inferior a la de
    cada capa por separado (797 de 800 payloads con ambas capas, frente a 794 de la librería con las
    cabeceras del proxy y 749 del WAF solo).

    \item \textbf{Montar el middleware después de todos los analizadores de cuerpo.}
    \textit{Acción:} registrar el middleware después de los analizadores de cuerpo de la aplicación,
    incluidos los de \texttt{multipart/form-data}. \textit{Justificación:} en las Rondas 1 a 3, la
    librería no inspeccionó 7 de 23 payloads de XSS enviados como \texttt{multipart/form-data}
    mientras el analizador se montaba por ruta, y bloqueó los 23 al montarlo de forma global antes
    del middleware. \textit{Resultado esperado:} que la librería inspeccione el contenido de todas
    las peticiones.

    \item \textbf{No desplegar la versión actual en rutas de alto volumen de datos.}
    \textit{Acción:} excluir de la protección, o paginar, las rutas que devuelven grandes volúmenes
    de datos por petición, hasta que se resuelva la saturación (recomendación~\ref{rec:cola}).
    \textit{Justificación:} en el escenario de volumen, el $\Delta p95$ fue de 6{,}495.46~ms y el
    7.00~\% de las peticiones agotó el tiempo de espera (Conclusión~\ref{conc:rendimiento}).
    \textit{Resultado esperado:} evitar una degradación de la disponibilidad del servidor anfitrión.
\end{enumerate}

\section{Mejoras técnicas de la librería}

\begin{enumerate}
    \setcounter{enumi}{4}
    \item \textbf{Usar el modo \texttt{monitor} por defecto.} \textit{Acción:} cambiar el modo por defecto del paquete, hoy \texttt{block}, a \texttt{monitor} en su próxima versión y documentar en el README cómo
    activar el bloqueo. \textit{Justificación:} las pruebas en aplicaciones de terceros se hicieron
    con la configuración por defecto y arrojaron bloqueos indebidos altos
    (Conclusión~\ref{conc:fp}). \textit{Resultado esperado:} que una instalación sin configuración
    no bloquee tráfico legítimo.

        \item \textbf{Verificar y, si se confirma, corregir el sesgo de la cabecera
    \texttt{User-Agent} y del recuento de cabeceras no estándar.} \textit{Acción:} reunir tráfico
    legítimo de navegadores y clientes HTTP diversos (incluidos cURL y los \texttt{User-Agent} por
    defecto de Node.js), evaluarlo con la configuración vigente, recalibrar el umbral de bloqueo con
    tráfico de varias aplicaciones, repetir las mediciones del Cuadro~\ref{tab:res-terceros-ua} y
    repetir la evaluación con la librería detrás de un proxy inverso, con y sin las cabeceras que este
    agrega. Si la hipótesis se confirma, la
    ampliación del corpus de entrenamiento corresponde al componente de Inteligencia Artificial y
    queda fuera de este documento. \textit{Justificación:} solo el 0.4~\% de las filas legítimas
    tiene un \texttt{User-Agent} largo, frente al 33.3--97.7~\% de las clases de ataque, y esa
    característica es la de mayor contribución del modelo (Conclusión~\ref{conc:fp}); la hipótesis
    no se verificó con un corpus legítimo alternativo. Además, cuatro cabeceras de proxy elevaron de
    0 a 4 el recuento de cabeceras fuera de la lista estándar y cambiaron el veredicto de 14 de los
    800 payloads en la Configuración~2 (Sección~\ref{sec:met-validez}); no se midió la importancia de
    esa característica en el modelo. \textit{Resultado esperado:} una reducción medible de los
    bloqueos indebidos o, si no se produce, el descarte de esa hipótesis, y la medida del efecto de
    las cabeceras del despliegue sobre el veredicto.

    \item \label{rec:cola}\textbf{Proteger al servidor ante saturación.} \textit{Acción:}
    instrumentar la cola de los hilos de inferencia durante el escenario de volumen y evaluar un
    mecanismo que deje pasar las peticiones (\textit{fail-open}) cuando la cola supere un límite.
    \textit{Justificación:} el 7.00~\% de tiempos de espera se concentró al final de cada
    repetición, un patrón consistente con una cola que crece, aunque la hipótesis no se comprobó
    (Conclusión~\ref{conc:rendimiento}). \textit{Resultado esperado:} confirmar o descartar la
    saturación y reducir los tiempos de espera del escenario de volumen.

    \item \textbf{Reducir el costo fijo de memoria y CPU.} \textit{Acción:} medir por separado el
    costo de cada componente de la librería (solo el Random Forest, el Random Forest con un hilo
    del Isolation Forest y la configuración actual) y evaluar \texttt{IF\_POOL\_SIZE=1} en lugar
    del valor actual de 2, repitiendo el benchmark con el mismo análisis pareado por repetición.
    \textit{Justificación:} el middleware agregó entre 331.9 y 375.9~MiB, más del doble del límite
    de 143.1~MiB, y alrededor de un núcleo de CPU, y ese costo no se atribuyó a ningún componente
    (Conclusión~\ref{conc:rendimiento}). \textit{Resultado esperado:} identificar el componente que
    domina el costo y la reducción que se logra con cada cambio de configuración.

    \item \textbf{Reducir la confusión entre categorías en los eventos registrados.}
    \textit{Acción:} revisar cómo se asigna la clase que se guarda en el evento y comprobarla con
    una matriz de confusión del registro. \textit{Justificación:} en la repetición de la
    Configuración~2 con el par~B, los 200 payloads de Command Injection se registraron como SQLi, y
    en DVNA, 27 de los 37 bloqueados de esa categoría se clasificaron como SQLi, aunque la decisión
    de bloqueo no cambió. \textit{Resultado esperado:} que el catálogo de ataques de la CLI sea
    confiable como herramienta de análisis.
\end{enumerate}

\section{Validaciones pendientes}

\begin{enumerate}
    \setcounter{enumi}{9}
    \item \textbf{Validar la librería en el entorno productivo de una startup guatemalteca.}
    \textit{Acción:} instalarla en modo \texttt{monitor} en una aplicación real y medir los
    bloqueos indebidos y el $\Delta p95$ con tráfico de usuarios. \textit{Justificación:} esta
    validación formaba parte del alcance complementario del protocolo y no se concretó por falta de
    tiempo para contactar a una startup durante el período de pruebas (Conclusión~\ref{conc:fp}).
    \textit{Resultado esperado:} medir los falsos positivos y el impacto en latencia con tráfico
    real y no solo con aplicaciones vulnerables de terceros.

    \item \textbf{Repetir las mediciones con el paquete vigente y con un corpus nuevo.}
    \textit{Acción:} con el paquete 0.1.1 del registro y el ejecutor corregido, repetir la cobertura con
    el paquete 0.1.0, las aplicaciones de terceros y el benchmark de rendimiento; repetir también
    las cuatro configuraciones de la Ronda~5 con un corpus nuevo, generado con otra semilla y un
    filtro de exclusión que incorpore, además de la igualdad exacta, la auditoría de subcadena
    normalizada, y fijar el digest de la imagen del WAF. Registrar el SHA-256
    del archivo publicado (el de la 0.1.1 es \texttt{abf8d5c4...}) y de sus modelos, y medir el
    desajuste entre la cadena de consulta con la que se entrenó y la que el middleware entrega al
    extractor. \textit{Justificación:} las cuatro configuraciones se midieron con la 0.1.1 sobre el corpus
    que reveló los defectos (en Path Traversal y Command Injection) y sobre un corpus de XSS
    construido después de conocer los resultados de la Ronda~5, y las aplicaciones de terceros y
    el benchmark no se midieron con ella (Conclusión~\ref{conc:cobertura}). \textit{Resultado
    esperado:} cifras independientes y atadas al archivo exacto evaluado, que confirmen o corrijan
    el 100.0~\% en Path Traversal, el 92.5~\% en Command Injection y el 100.0~\% en XSS, y una
    comparación con el WAF sobre un corpus independiente.

        \item \textbf{Ampliar la explotabilidad confirmada de XSS más allá de un caso.}
    \textit{Acción:} confirmar con OWASP ZAP u otras herramientas, sobre más endpoints o
    aplicaciones, la explotabilidad de payloads del corpus propio, en particular los dirigidos a
    \texttt{content} y \texttt{title}, cuya vulnerabilidad no se confirmó, y repetir el filtrado
    contra el entrenamiento con el criterio de subcadena normalizada. \textit{Justificación:} la
    cobertura de XSS se midió sobre 200 payloads (200 de 200), pero la explotabilidad confirmada
    descansa en un único caso (1 de 1, ZAP) y 136 de los 200 payloads apuntan a campos sin
    vulnerabilidad confirmada (Conclusión~\ref{conc:cobertura}). \textit{Resultado esperado:} una
    medición de explotabilidad de XSS con más de un caso confirmado, que permita separar la
    detección de la explotabilidad.

    \item \textbf{Investigar los payloads de Command Injection que la combinación no bloquea.}
    \textit{Acción:} analizar por qué \texttt{cmdi\_6}, \texttt{cmdi\_184} y \texttt{cmdi\_185} no
    los bloquea ninguna configuración, y evaluar si el veredicto de la librería debe depender de
    las cabeceras que agrega la infraestructura, repitiendo las mediciones con y sin ellas.
    \textit{Justificación:} en la medición del 25 de octubre de 2026 esos tres son los únicos
    payloads que la combinación no bloquea (el WAF no activa ninguna regla sobre ellos, y con las
    cabeceras del proxy la probabilidad de la clase de ataque más probable del Random Forest queda
    en 0.3333 o menos), y cuatro payloads de la misma categoría solo se bloquean cuando el proxy
    agrega cuatro cabeceras \texttt{X-*} (Conclusión~\ref{conc:waf}). \textit{Resultado esperado:}
    saber si los tres payloads son una evasión real de las dos capas y cuánto cambia el veredicto
    de la librería según las cabeceras del despliegue.

    \item \textbf{Monitorear la memoria en un escenario sostenido de condición única.}
    \textit{Acción:} ejecutar un escenario de 30~minutos o más con una sola condición estable.
    \textit{Justificación:} la corrida continua de 41.9~minutos mezcla condiciones y no permite
    descartar una fuga de memoria dentro de un proceso estable (Sección~\ref{sec:met-restricciones}).
    \textit{Resultado esperado:} confirmar o descartar una fuga de memoria.

    \item \textbf{Medir la contribución del Isolation Forest con las pruebas de penetración.}
    \textit{Acción:} registrar en cada ronda de pruebas el veredicto \texttt{pass\_anomaly} de cada
    payload que el Random Forest no bloquea y de cada petición legítima, y reportar cuántos
    ataques no bloqueados quedan marcados como anómalos y cuántas alertas generan las peticiones
    legítimas. \textit{Justificación:} la contribución del detector de anomalías a la detección de
    variantes solo tiene evidencia indirecta: en un conjunto de validación ajeno a este trabajo marca
    19 de los 39 ataques que el Random Forest deja pasar, pero bloquear sus marcas agregaría 866
    benignos bloqueados \cite{valenzuela2026ia} (Conclusión~\ref{conc:variantes}).
    \textit{Resultado esperado:} estimar, con el tráfico de ataque de la evaluación del sistema,
    cuánto aporta el detector y cuántas alertas adicionales genera.
\end{enumerate}

\section{Reproducibilidad y entrega de datos}

\begin{enumerate}
    \setcounter{enumi}{15}
    \item \textbf{Publicar los vectores derivados y depurar el conjunto de pruebas extremo a
    extremo.} \textit{Acción:} publicar los vectores de características derivados como recurso
    descargable de una versión del repositorio, sin la columna \texttt{\_row\_hash} y después de
    revisar que ninguna característica permita reconstruir el texto de la petición; y sustituir, en
    el conjunto de pruebas extremo a extremo versionado, los payloads cuyo texto proviene de
    fuentes sin licencia especificada. \textit{Justificación:} el Alcance entrega los scripts y los
    vectores derivados, y las fuentes \texttt{XSS\_dataset.csv} y \texttt{command injection.csv}
    (3.4~\% de las filas) no tienen licencia especificada por su origen
    (Sección~\ref{subsec:met-f1}). \textit{Resultado esperado:} un repositorio público coherente
    con lo que afirma la auditoría de fuentes.

    \item \textbf{Versionar el generador de los conjuntos sintéticos de Command Injection.}
    \textit{Acción:} incorporar al repositorio el script que los genera, o declarar esa porción del
    corpus como no reproducible. \textit{Justificación:} el 17.8~\% de las filas de Command
    Injection del conjunto de prueba (261 de 1{,}468) proviene de esos conjuntos y su generador no
    está versionado (Conclusión~\ref{conc:trazabilidad}). \textit{Resultado esperado:}
    reproducibilidad completa de la partición de prueba.
\end{enumerate}

\section{Ampliación del alcance}

\begin{enumerate}
    \setcounter{enumi}{17}
    \item \textbf{Desacoplar la detección de Node.js y de npm.} \textit{Acción:} separar un
    núcleo de detección reutilizable y escribir adaptadores para otros ecosistemas (por ejemplo,
    \textit{frameworks} de Python o de Java, o un \textit{proxy} local), verificando en cada uno la
    paridad de las predicciones con la de referencia, como se hizo entre Python y ONNX.
    \textit{Justificación:} el Alcance excluyó otros entornos de ejecución, pero el extractor, los
    modelos en ONNX y el umbral no dependen del lenguaje; lo que está atado a Node.js es el
    adaptador de Express y la ejecución en hilos de trabajo (Sección~\ref{subsec:fuera-alcance}).
    \textit{Resultado esperado:} ampliar la población de aplicaciones que pueden usar la detección
    sin reescribir el modelo.

    \item \textbf{Exponer la librería como servidor del Protocolo de Contexto de Modelo (MCP).}
    \textit{Acción:} publicar un servidor MCP que permita a un asistente de programación consultar
    el catálogo de ataques y los \textit{endpoints} más atacados que ya ofrece la CLI, definiendo
    antes qué operaciones son de solo lectura y cuáles podrían cambiar el modo o el umbral.
    \textit{Justificación:} la CLI ya expone esa información, y un asistente con acceso de
    escritura ampliaría la superficie de ataque de la propia herramienta de protección.
    \textit{Resultado esperado:} consulta asistida de los eventos sin ampliar esa superficie.
\end{enumerate}

\section{Para trabajos de investigación similares}

\begin{enumerate}
    \setcounter{enumi}{19}
    \item \textbf{Definir un conjunto de evaluación externo desde el inicio.} \textit{Acción:}
    reunir tráfico legítimo de varias aplicaciones y no limitarse a una partición de la misma
    distribución que el entrenamiento. \textit{Justificación:} las métricas del corpus (F1 de
    0.9239 a 0.9956) no predijeron los bloqueos con tráfico legítimo real (Conclusión~\ref{conc:fp}).
    \textit{Resultado esperado:} métricas que reflejen el comportamiento en producción y no solo la
    separación entre clases dentro del corpus.

    \item \textbf{Medir y tratar los casi-duplicados antes de particionar.} \textit{Acción:}
    medirlos (por ejemplo, con MinHash), agrupar por plantilla al dividir en entrenamiento,
    validación y prueba, y repetir la evaluación con el corpus deduplicado. \textit{Justificación:}
    SQLi tiene unas 395 veces más pares casi duplicados que filas y es el 59.3~\% del conjunto de
    prueba, un efecto que este trabajo no cuantificó (Conclusión~\ref{conc:clasificacion}).
    \textit{Resultado esperado:} cuantificar cuánto de las métricas actuales depende de esa
    saturación.

    \item \textbf{Bloquear el conjunto de prueba una sola vez y versionar cada medición.}
    \textit{Acción:} bloquear el conjunto de prueba una vez para todo el proyecto, atar cada cifra
    al hash de los datos y del modelo que la produjeron, y guardar cada medición de rendimiento
    como un artefacto versionado con el commit, el hardware, la versión del entorno de ejecución y
    las muestras crudas separadas por condición, como se hizo con el benchmark del extractor y con
    la evidencia de la corrida de latencia. \textit{Justificación:} el valor del bloqueo cambió
    nueve veces y el conjunto de prueba se leyó tres veces en la generación vigente
    (Conclusión~\ref{conc:trazabilidad}). \textit{Resultado esperado:} cifras atables a un único
    conjunto de datos y a un único modelo.

    \item \textbf{Construir los corpus de evaluación con un generador versionado y un filtro de
    exclusión.} \textit{Acción:} versionar el generador y aplicar el filtro contra el
    entrenamiento desde la primera ronda, y reportar siempre la tasa de detección junto con la de
    falsos positivos. \textit{Justificación:} la Ronda~4 tuvo fuga confirmada en Command Injection,
    que redujo la cobertura atribuible de 100~\% a 91.5~\% al corregirse
    (Conclusión~\ref{conc:trazabilidad}). \textit{Resultado esperado:} cifras de detección que no
    estén infladas por solapamiento con el entrenamiento.
\end{enumerate}

\ifdefined\modoensamble
\else
	\section*{Referencias}
	\printbibliography[heading=none]
	\end{document}
\fi
